Evropské systémy sociálního dialogu se~liší v~tom, jakým způsobem převádějí vyjednávací sílu sociálních partnerů do~plošných mzdových a~pracovních standardů. Komparativní výzkum rozlišuje několik vyhraněných modelů, které pro strukturu českého trhu práce představují klíčové kontrastní body. Jedná se zejména o~dánský model \textit{flexicurity} propojující vysokou flexibilitu zaměstnavatelů s~masivními výdaji na~\aclacc{APZ} a~univerzálním pokrytím kolektivními smlouvami, rakouský komorový systém (\textit{Kammer-system}) zajišťující plošný dosah dohod i~při nižší odborové organizovanosti a~německý model spolurozhodování (\textit{Mitbestimmung}) založený na odvětvových standardech a~podnikové participaci. V~kontrastu s~těmito koordinovanými systémy stojí decentralizované trajektorie v~Polsku a~na~Slovensku, jež jsou svou závislostí na~podnikovém vyjednávání a~nízkým smluvním pokrytím blíže situaci v~\acgen{ČR}.~\cite{Hala_SystemySocialnihoDialogu2013}\par
